% !TeX program = lualatex
% Compile twice: lualatex 130.tex
% All references and prose are embedded; no BibTeX database or external inputs.
\documentclass[11pt,a4paper]{article}
\usepackage[margin=24mm,headheight=14pt]{geometry}
\usepackage{fontspec}
\setmainfont{Latin Modern Roman}
\setsansfont{Latin Modern Sans}
\setmonofont{Latin Modern Mono}
\usepackage{amsmath,amssymb,mathtools}
\usepackage{unicode-math}
\setmathfont{Latin Modern Math}
\usepackage{microtype}
\usepackage{enumitem,xurl,needspace}
\usepackage[dvipsnames]{xcolor}
\usepackage{hyperref}
\hypersetup{unicode=true,colorlinks=true,linkcolor=MidnightBlue,urlcolor=MidnightBlue,
 pdftitle={500 Mathematical Problem Statements},pdfauthor={Source-annotated compilation},bookmarksnumbered=true}
\usepackage{bookmark}
\usepackage{tocloft}
\setlength{\cftsecnumwidth}{3em}
\cftsetpnumwidth{2.5em}
\cftsetrmarg{4em}
\usepackage{titlesec}
\titleformat{\section}{\Large\bfseries\raggedright}{\thesection}{0.7em}{}
\usepackage{fancyhdr}
\pagestyle{fancy}\fancyhf{}
\fancyhead[L]{\small\sffamily Mathematical problem statements}
\fancyhead[R]{\small Source edition 22}
\fancyfoot[C]{\thepage}
\setlength{\parindent}{0pt}
\setlength{\parskip}{5pt plus 1pt}
\setlength{\emergencystretch}{3em}
\setcounter{tocdepth}{1}
\setcounter{secnumdepth}{1}
\setlist[itemize]{leftmargin=3em,itemsep=5pt,topsep=4pt,parsep=0pt}
\allowdisplaybreaks
\newcommand{\N}{\mathbb{N}}
\newcommand{\Z}{\mathbb{Z}}
\newcommand{\Q}{\mathbb{Q}}
\newcommand{\R}{\mathbb{R}}
\newcommand{\C}{\mathbb{C}}
\newcommand{\F}{\mathbb{F}}
\newcommand{\A}{\mathbb{A}}
\newcommand{\PP}{\mathbb P}
\newcommand{\E}{\mathbb{E}}
\newcommand{\eps}{\varepsilon}
\newcommand{\dd}{\,\mathrm d}
\newcommand{\sref}[2]{\hyperlink{source:#1:#2}{[S#2]}}
\newcommand{\eref}[2]{\hyperlink{extra:#1:#2}{[E#2]}}
% Turn the next switch off for a shorter reading copy. The quotations preserve
% every exactTarget field, including qualifications omitted from a short summary.
\newif\ifshowcatalogue
\showcataloguetrue
\newcommand{\cataloguescope}[1]{\ifshowcatalogue
 \paragraph{Exact scope in the supplied catalogue.}
 \begin{quote}\small #1\end{quote}\fi}
\hypersetup{pdftitle={Problem 130 - Reciprocal sums and progressions},pdfauthor={Open Problems Math research notebook}}
\fancyhead[L]{\small\sffamily Open Problems Math \textbar\ Research notebook}
\fancyhead[R]{\small\sffamily 130 \textbar\ 27 September 2026}
\numberwithin{equation}{section}
\begin{document}
\setcounter{section}{129}
% ================================================================
% problemId: problem.erdos-problem-3-divergent-harmonic-series-and-long-arithmetic-progressions
% source releaseRank: 130; source releaseStatus: open_with_solved_subcases
% exactTarget SHA-256: 38149aeefa1dff68afa64847857b023ebf1c0a2a076d58b0ea361a371b74019d
\clearpage
\section{\texorpdfstring{Erdős problem 3: divergent harmonic series and long arithmetic progressions}{Erdős problem 3: divergent harmonic series and long arithmetic progressions}}\label{130}
\subsection{Definitions and mathematical statement}
For $A\subseteq\Z_{>0}$, divergence of its reciprocal sum means $\sum_{a\in A}1/a=\infty$. A nonconstant arithmetic progression of length $k$ is $a,a+d,\ldots,a+(k-1)d$ with integers $a,d\ge1$. The question is whether
\[
 \sum_{a\in A}\frac1a=\infty\quad\Longrightarrow\quad
 \forall k\ge1\ \exists a,d\ge1:\ \{a+jd:0\le j<k\}\subseteq A.
\]
Positive density is not assumed, and the common difference can depend on the requested length.
\par\smallskip\textit{Formulation and concept references:} \sref{130}{1}.
\cataloguescope{Determine whether every subset A of the natural numbers whose reciprocal sum diverges must contain arbitrarily long arithmetic progressions.}
\subsection{Short English statement}
Does divergence of a set's reciprocal sum force it to contain arithmetic progressions of every finite length?
% BEGIN RESEARCH ADDITION 130
\subsection{Research attempt}

\paragraph{Current frontier and scope.}
The reciprocal-sum question is much stronger than a positive-density
progression theorem: the density of a set with divergent reciprocal sum
may tend to zero.  The supplied entry associates the problem with
Green's historical article. \sref{130}{1} The exact target used here is
the displayed catalogue statement; the quantitative input used below
is the separate theorem of Bloom--Sisask.
For length three, Bloom--Sisask's Theorem 1 gives
$r_3(N)\leq N\exp(-c(\log N)^{1/9})$, which is strong enough to make
the reciprocal sum of any three-term-progression-free set finite.
\eref{130}{1}  This does not supply the same quantitative conclusion
for every fixed progression length.  The frontier check on
27 September 2026 did not yield a verified resolution of the full
all-length statement.  The deductions below are independent of any
purported solution claim.

\paragraph{Attack map.}
A density-increment approach must produce a summable sequence of density
upper bounds across scales, not merely densities tending to zero.
A counterexample approach might glue large finite progression-free sets
at separated scales; the difficulty is preventing progressions that cross
between blocks.  A compactness argument based only on finite examples
would not control their reciprocal mass.  Here the gluing route can be
made exact, producing an equivalence between the original fixed-length
problem and a specific numerical series, together with a canonical
computable counterexample if that series diverges.

\paragraph{Attempt I: separate blocks so that no three-term progression crosses.}
Fix $k\geq3$.  Write
\[
 r_k(N)=\max_{\substack{B\subseteq\{1,\ldots,N\}\\
            B\text{ has no nonconstant }k\text{-term progression}}}|B|,
 \qquad a_j=\frac{r_k(4^j)}{4^j}\quad(j\geq0).
\]
For each $j$, choose a largest $k$-term-progression-free set
$B_j\subseteq[1,4^j]\cap\Z$.  Define
\begin{equation}\label{130:gluing}
 A_k=\bigcup_{j\geq0}(2\cdot4^j+B_j),\qquad
 I_j=[2\cdot4^j+1,3\cdot4^j]\cap\Z.
\end{equation}
The intervals $I_j$ are disjoint.  The separation is chosen so that
\emph{no} nonconstant three-term progression in their union can meet
two different intervals.

To prove the latter assertion, suppose $x<y<z$ and $x+z=2y$, with
$z\in I_j$ the highest occupied block.  If $y$ lies in an earlier
block, then $y\leq3\cdot4^{j-1}$ and
\[
 z=2y-x\leq6\cdot4^{j-1}-1<2\cdot4^j+1,
\]
a contradiction.  If $y\in I_j$ but $x$ lies earlier, then
\[
 z=2y-x\geq4\cdot4^j+2-3\cdot4^{j-1}>3\cdot4^j,
\]
again a contradiction.  For $j=0$ there is no earlier block.  Thus
all three terms lie in one interval.

Every consecutive triple of a $k$-term progression is a three-term
progression.  Overlapping consecutive triples force all $k$ terms
into the same block.  Such a progression cannot occur in $A_k$, since
its translate back into $B_j$ is forbidden.  This proves that $A_k$
is $k$-term-progression-free even when the individual $B_j$ contain
many shorter progressions.  Assuming they are three-term-free would
have unnecessarily weakened the construction.

\paragraph{Attempt II: an exact summability equivalence.}
Every member of $2\cdot4^j+B_j$ is at most $3\cdot4^j$.  Consequently
\begin{equation}\label{130:lower}
 \sum_{a\in A_k}\frac1a
 \geq\frac13\sum_{j\geq0}\frac{|B_j|}{4^j}
 =\frac13\sum_{j\geq0}a_j.
\end{equation}
Conversely, for any $k$-term-progression-free $A\subseteq\Z_{>0}$,
partition the positive integers into $[4^j,4^{j+1})$.  There are at
most $r_k(4^{j+1})$ members of $A$ in this interval, and each has
reciprocal at most $4^{-j}$.  Therefore
\begin{equation}\label{130:upper}
 \sum_{a\in A}\frac1a\leq4\sum_{j\geq1}a_j.
\end{equation}
Both statements use nonnegative series, so no conditionally convergent
rearrangement is involved.

\textbf{Lacunary extremal-series lemma.}
For every fixed integer $k\geq3$, the following are equivalent:
\[
 \begin{gathered}
 \text{every $k$-term-progression-free subset of $\Z_{>0}$}\\
 \text{has finite reciprocal sum}\\[2pt]
 \Longleftrightarrow\quad
 \sum_{j\geq0}\frac{r_k(4^j)}{4^j}<\infty.
 \end{gathered}
\]
\emph{Proof.} Convergence implies the first assertion by
\eqref{130:upper}.  Divergence produces the explicit counterexample
\eqref{130:gluing} by \eqref{130:lower}. \hfill$\square$

It follows that the entire conjecture in the original entry is
\emph{equivalent} to convergence of this series for every fixed
$k\geq3$.  Lengths one and two cause no exception: a set with divergent
reciprocal sum is infinite.  This is more than a one-directional
sufficient density estimate, but no assertion is made that the
series has now been evaluated for general $k$.

The same proof gives a quantitative finite form.  Put
\[
 W_k(N)=\max_{\substack{A\subseteq[1,N]\cap\Z\\A\text{ $k$-term-free}}}
          \sum_{a\in A}\frac1a.
\]
For every $J\geq0$,
\begin{equation}\label{130:finite}
 \frac13\sum_{j=0}^{J}a_j
 \leq W_k(3\cdot4^J)
 \leq4\sum_{j=1}^{J+1}a_j.
\end{equation}
The lower bound uses the first $J+1$ blocks in \eqref{130:gluing}; the
upper bound is the same interval partition, stopped before $4^{J+1}$.
Thus convergence would in fact give a \emph{uniform} reciprocal-sum
bound depending only on $k$, although the original statement does not
ask for an explicit bound or a rate.

\paragraph{Attempt III: make a hypothetical counterexample computable.}
Choose $B_j$ canonically: enumerate all subsets of $[1,4^j]$, test
all candidate progressions, maximize cardinality, and use lexicographic
order to break ties.  This is an astronomically slow but terminating
finite procedure.  Membership of an integer $n$ in $A_k$ is decidable:
find whether it belongs to one of the disjoint intervals $I_j$, and,
if so, perform that finite search and test $n-2\cdot4^j\in B_j$.

Therefore, if the conjecture is false at some fixed length $k$, it
has a \emph{computable} counterexample of the specific form
\eqref{130:gluing}.  No effectiveness assumption has been appended
to the original problem: the conclusion follows from the extremal
series equivalence.  The unknown point is not the algorithm's
termination, but whether its output set has divergent reciprocal sum.
Computing larger finite portions cannot decide that infinite-series
question without a new uniform argument.

A nonconstructive choice of different near-extremizers at every scale
is likewise not an obstacle here.  The separation lemma permits all
scales to be used simultaneously.  The missing arithmetic information
is entirely in their attainable relative sizes $a_j$.

\paragraph{Testing the reduction against known estimates.}
The length-three theorem cited above yields
\[
 a_j\leq\exp(-c'(j\log4)^{1/9})
\]
after adjusting constants and omitting finitely many initial terms.
This series converges, for example by comparison with the integral
of $\exp(-c''x^{1/9})$.  Hence the three-term case follows rigorously
from the extremal-series lemma. \eref{130}{1}
No comparable all-$k$ summable estimate has been proved in this attempt.

For general $k$, a bound
$r_k(N)\ll_k N/(\log N)^{1+\delta_k}$ with any $\delta_k>0$ would
suffice.  In contrast, $r_k(N)=o(N)$ gives only $a_j\to0$, which does
not imply $\sum a_j<\infty$.  Even an upper bound of order $1/j$
for $a_j$ would not prove convergence.  These observations describe
what particular estimates can or cannot establish; they do not assert
that the actual extremal sequence equals $1/j$ or that a counterexample
has been found.

The harmonic weights are essential to the calculation: translation to
$[2\cdot4^j+1,3\cdot4^j]$ keeps the reciprocal mass comparable to
$|B_j|/4^j$.  Moving blocks arbitrarily far out would eliminate crossing
progressions but could destroy all the reciprocal divergence.

\paragraph{Stress test.}
The gluing proof is checked for the worst placement of the smaller
terms, not merely for progressions with one term in each of three
successive blocks.  It handles two terms in the largest block and one
term in any earlier block, and also one term in the largest block.
Only nonconstant progressions are considered, exactly as required.
The overlapping-triple argument is what extends the construction from
length three to every $k\geq3$.

The companion script exhaustively computes small extremizers for
$k=3,4$ and $N\leq16$, forms the corresponding initial blocks, and
checks the absence of cross-block three-term progressions and of
$k$-term progressions in their union.  Reciprocal sums and both sides
of the finite inequalities are checked with exact rational arithmetic.
This verifies examples and endpoints; the all-scale proof is supplied
above, rather than inferred from the finite search.

The exact search, with lexicographic tie-breaking, gives the following
initial extremizers used in the finite gluing check:
\[
 \begin{array}{c|r|r|l}
 k&N&r_k(N)&B\\\hline
 3&1&1&\{1\}\\
 3&4&3&\{1,2,4\}\\
 3&16&8&\{1,2,4,5,10,11,13,14\}\\
 4&1&1&\{1\}\\
 4&4&3&\{1,2,3\}\\
 4&16&10&\{1,2,3,5,6,8,9,10,15,16\}\\
 \end{array}
\]
The last column records actual subsets, not just a numerical upper
bound; maximality in this finite range is certified by exhaustive search.

\paragraph{Outcome.}
\textbf{Outcome: Conditional reduction.}

The fixed-length reciprocal-sum assertion is exactly equivalent to
summability of $r_k(4^j)/4^j$, with a canonical computable counterexample
if the series diverges.  The construction also gives uniform finite
harmonic-mass inequalities and recovers the length-three case from an
established quantitative theorem.  For arbitrary $k$, neither
convergence nor divergence of the required series has been established.
These reductions are not claimed to be new to the literature, and no
complete solution of the original all-length problem is claimed.
% END RESEARCH ADDITION 130
\subsection{Sources}
\begin{itemize}\raggedright
\item[\textnormal{[S1]}]\hypertarget{source:130:1}{} B. Green, Arithmetic progressions at the Journal of the LMS, Journal of the London Mathematical Society (2026).
\url{https://londmathsoc.onlinelibrary.wiley.com/doi/abs/10.1112/jlms.70483}.
{\small\textit{Catalogue formulation source.}}
% BEGIN RESEARCH REFERENCES 130
\item[\textnormal{[E1]}]\hypertarget{extra:130:1}{} Thomas F. Bloom and Olof Sisask.
\emph{An improvement to the Kelley--Meka bounds on three-term arithmetic progressions}, arXiv:2309.02353v1, 5 September 2023. Theorem 1: $r_3(N)\leq N\exp(-c(\log N)^{1/9})$. Used as a sufficient summable bound, not as an all-length theorem.
\url{https://arxiv.org/html/2309.02353v1}.
{\small\textit{Consulted 27 September 2026 for the indicated result or scope.}}
% END RESEARCH REFERENCES 130
\end{itemize}


\end{document}
